\SourceLeaf{345}
\par\medskip\Needspace{4\baselineskip}\pdfbookmark[2]{Art. 7. 8.}{nachlass-biquadratic-III-art-7-8}\centerline{\textbf{7. 8.}}\nobreak\smallskip
Diese Regel ist allgemein, was für eine Zahl auch $M$ bedeute. Für den Fall, der zunächst den Gegenstand unserer Untersuchung ausmachen soll, wo $M$ ungerade und von der Form $1+(2+2i)N$ vorausgesetzt wird, ist eine etwas abgeänderte Vorschrift zweckmässiger.

Man denke sich die Zahlen $f$ wiederum in $4$ Classen zerlegt; in die erste setzt man die $(h)$, deren Doppeltes sich auch noch in $f$ findet; in die zweite $h'$ zählen wir die, deren Doppelte $2h'$ zu $f'$ gehören, und ebenso $h''$ und $h'''$ bedeuten diejenigen, deren Doppelte zu $f''$ und $f'''$ gehören. Es ist also der Decident $\varepsilon$
\[
\varepsilon=\Sigma\Theta\frac{2hM}{m}+\Sigma\Theta\frac{2h'M}{m}+\Sigma\Theta\frac{2h''M}{m}+\Sigma\Theta\frac{2h'''M}{m}
\]
Den Complexus aller $2h$ und $-2h''+(1+i)m$ nennen wir $H$

den von allen $-i(2h'-m)$ und $i(2h'''-im)$ nennen wir $H'$

$H$ und $H'$ umfassen also alle $f$, jene sind die geraden, diese die ungeraden.

Ferner sind folgende Relationen in Anwendung zu bringen
\[
\begin{aligned}
\Theta iN&=1+\Theta N\\
\Theta(-N)&=2+\Theta N\\
\Theta(-iN)&=3+\Theta N\\
\Theta(N+1)&=1-\Theta N\\
\Theta(N+1+i)&=2+\Theta N\\
\Theta(N+i)&=3-\Theta N
\end{aligned}
\]
folglich
\[
\begin{aligned}
\Theta\frac{(-2h'i+mi)M}{m}&=3-\Theta\frac{-2h'iM}{m}=-\Theta\frac{2h'M}{m}\\
\Theta\frac{(-2h''+m(1+i))M}{m}&=2+\Theta\frac{-2h''M}{m}=\Theta\frac{2h''M}{m}\\
\Theta\frac{(2h'''i+m)M}{m}&=1-\Theta\frac{2h'''iM}{m}=-\Theta\frac{2h'''M}{m}
\end{aligned}
\]
und
\[
\begin{aligned}
\varepsilon={}&\Sigma\Theta\frac{2hM}{m}-\Theta\frac{(-2h'i+mi)M}{m}\\
&+\Theta\frac{(-2h''+m(1+i))M}{m}-\Theta\frac{(2h'''i+m)M}{m}\\
={}&\Sigma\Theta\frac{HM}{m}-\Sigma\Theta\frac{H'M}{m}
\end{aligned}
\]

\SourceLeaf{346}
oder
\[
\varepsilon=\Sigma\pm\Theta\frac{fM}{m}
\]
ubi signum superius accipiendum pro paribus $f$, inferius pro imparibus.

\Article{9}
Es sei nun allgemein $f=\xi+\eta i$. Die Zahlen $\xi,\eta$ sind durch die Bedingung, dass $f$ ein kleinster Rest von $m$ sein, oder $\frac{f}{m}=x+yi$ gesetzt, $x$ und $y$ zwischen den Grenzen $0$ und $\frac12$ liegen müssen, innerhalb gewisser Grenzen beschränkt, wofür sich durch Unterscheidung der verschiedenen Fälle leicht bestimmte Regeln geben liessen. Ertheilen wir $\eta$ einen bestimmten Werth, so wird wiederum $\xi$ seine bestimmten Grenzen haben. Z.\,B. wenn wir annehmen, dass $a$ negativ, $b$ positiv ist, so muss, da
\[
\begin{aligned}
x&=\frac{a\xi+b\eta}{aa+bb}\\
y&=\frac{a\eta-b\xi}{aa+bb}
\end{aligned}
\]
\[
\begin{array}{rll}
\mathrm{I.}&\text{damit }x\text{ positiv werde}&\xi<-\frac{b}{a}\eta\\[3pt]
\mathrm{II.}&\text{damit }y\text{ positiv werde}&\xi<\frac{a}{b}\eta\\[3pt]
\mathrm{III.}&\text{damit }x<\frac12\text{ werde}&\xi>\frac{aa+bb-2b\eta}{2a}\\[3pt]
\mathrm{IV.}&\text{damit }y<\frac12\text{ werde}&\xi>\frac{2a\eta-aa-bb}{2b}
\end{array}
\]

für positive $\eta$ schliesst die zweite Bedingung bereits die erste ein, für negative $\eta$ hingegen ist es umgekehrt; ebenso ist die dritte Bedingung schon in der vierten enthalten,
\[
\begin{array}{ll}
\text{wenn}&\eta<\tfrac12(a+b)\\
\text{und umgekehrt, wenn}&\eta>\tfrac12(a+b)
\end{array}
\]

Wir haben indessen nicht nöthig alle acht Fälle, die hier eintreten können, besonders zu betrachten, sondern bezeichnen nur für einen bestimmten Werth von $\eta$ die kleinere Grenze von $\xi$ durch $\xi^0$, die grössere durch $\xi^{00}$ und bemerken nur, dass bei diesen Grenzwerthen immer entweder $x=0,y=0,x=\frac12,y=\frac12$ ist, und zwar dass
\SourceLeaf{347}
\[
\begin{array}{l|l|l}
\text{wenn}&\text{in der \emph{obern} Grenze}&\text{in der \emph{untern} Grenze}\\
a\text{ pos. }b\text{ positiv}&x=\frac12\text{ oder }y=0&x=0\text{ oder }y=\frac12\\
a\text{ neg. }b\text{ positiv}&x=0\text{ oder }y=0&x=\frac12\text{ oder }y=\frac12\\
a\text{ neg. }b\text{ negativ}&x=0\text{ oder }y=\frac12&x=\frac12\text{ oder }y=0\\
a\text{ pos. }b\text{ negativ}&x=\frac12\text{ oder }y=\frac12&x=0\text{ oder }y=0
\end{array}
\]
sein muss. Wir werden diese vier Fälle Kürze halber so unterscheiden, dass wir sagen, im ersten gehöre $m$ zum ersten Quadranten, im zweiten zum zweiten etc.

\Article{10}
Wir wollen nun das Aggregat aller $\pm\Theta\frac{fM}{m}$ näher betrachten, bei denen $\eta$ einen bestimmten Werth hat. Indem $\xi$ nach und nach stetig von dem kleinsten Werthe $\xi^0$ bis zum grössten $\xi^{00}$ wächst, wird sich
\[
\frac{(\xi+\eta i)M}{m}=X+Yi
\]
auch nach dem Gesetze der Stetigkeit ändern, und zwar wird, wenn $\frac{M}{m}$ im ersten Quadranten liegt, sowohl $X$ als $Y$ beständig wachsen; liegt $\frac{M}{m}$ im zweiten Quadranten, so wird $X$ beständig abnehmen und $Y$ zunehmen; im dritten Quadranten wird das umgekehrte vom ersten, im vierten das umgekehrte vom zweiten Statt finden. Allein die in $X+iY$ enthaltene ganze Zahl wird sich sprungsweise ändern, indem entweder $[X]$ oder $[Y]$ sich um Eine Einheit ändert. Es seien die Werthe von $\xi$, wo ein solcher Uebergang Statt findet, d.\,i. wo entweder $X$ oder $Y$ eine ganze Zahl wird, der Reihe nach folgende
\[
\xi',\ \xi'',\ \xi'''\dotsfour\xi^n
\]

Hier muss bemerkt werden, dass weder diese Werthe noch $\xi^0$ und $\xi^{00}$ ganze Zahlen sein können, ausgenommen für $\eta=0$, wo entweder $\xi^0$ oder $\xi^{00}=0$ wird. Es sei nun
\[
\begin{aligned}
\Theta\frac{(\xi'+\eta i)M}{m}-\Theta\frac{(\xi^0+\eta i)M}{m}&=\delta'\quad\text{(anders auszudrücken)}\\
\Theta\frac{(\xi''+\eta i)M}{m}-\Theta\frac{(\xi'+\eta i)M}{m}&=\delta''\\
&\text{etc.}\\
\Theta\frac{(\xi^{00}+\eta i)M}{m}-\Theta\frac{(\xi^{n-1}+\eta i)M}{m}&=\delta^n
\end{aligned}
\]
\SourceLeaf{348}
so sieht man leicht, weil zwischen $\xi^0$ und $\xi'$ $[\frac12\xi']-[\frac12\xi^0]$ gerade und $[\frac12\xi'+\frac12]-[\frac12\xi^0+\frac12]$ ungerade ganze Zahlen liegen etc., dass, blos den bestimmten Werth von $\eta$ betrachtet,
\[
\begin{aligned}
(\pm1)\Sigma\Theta\frac{fM}{m}
={}&\{[\tfrac12\xi']-[\tfrac12\xi^0]-[\tfrac12\xi'+\tfrac12]+[\tfrac12\xi^0+\tfrac12]\}\\
&\qquad\pdot\Theta\frac{(\xi^0+\eta i)M}{m}\\
&+\{[\tfrac12\xi'']-[\tfrac12\xi']-[\tfrac12\xi''+\tfrac12]+[\tfrac12\xi'+\tfrac12]\}\\
&\qquad\pdot\{\Theta\frac{(\xi^0+\eta i)M}{m}+\delta'\}\\
&+\{[\tfrac12\xi''']-[\tfrac12\xi'']-[\tfrac12\xi'''+\tfrac12]+[\tfrac12\xi''+\tfrac12]\}\\
&\qquad\pdot\{\Theta\frac{(\xi^0+\eta i)M}{m}+\delta'+\delta''\}\\
&+\text{etc.}\\
&+\{[\tfrac12\xi^{00}]-[\tfrac12\xi^n]-[\tfrac12\xi^{00}+\tfrac12]+[\tfrac12\xi^n+\tfrac12]\}\\
&\qquad\pdot\{\Theta\frac{(\xi^0+\eta i)M}{m}+\delta'+\delta''+\dotstwo+\delta^n\}
\end{aligned}
\]
\[
\begin{aligned}
={}&-([\tfrac12\xi^0]-[\tfrac12\xi^0+\tfrac12])\pdot\Theta\frac{(\xi^0+\eta i)M}{m}\\
&-([\tfrac12\xi']-[\tfrac12\xi'+\tfrac12])\pdot\delta'\\
&-([\tfrac12\xi'']-[\tfrac12\xi''+\tfrac12])\pdot\delta''\\
&-\text{etc.}\\
&-([\tfrac12\xi^n]-[\tfrac12\xi^n+\tfrac12])\pdot\delta^n\\
&+([\tfrac12\xi^{00}]-[\tfrac12\xi^{00}+\tfrac12])\pdot\Theta\frac{(\xi^{00}+\eta i)M}{m}
\end{aligned}
\]
(wo das obere Zeichen für gerade $\eta$, das untere für ungerade gilt.)

Die Zahlen $\delta',\delta'',\delta'''$ u.\,s.\,w. können keine andere Werthe haben als $+1$ und $-1$. Den Werth $+1$ bekommt $\delta'$, wenn, die Werthe von $X,Y$, die zu $\xi'$ gehören, durch $X',Y'$ bezeichnet,
\[
\begin{array}{l|l}
\frac{M}{m}\text{ im 1. Quadr.}&\frac{M}{m}\text{ im 2. Quadr.}\\
X'\text{ ganze gerade Zahl}&X'\text{ ganze gerade Zahl}\\
\text{und }[Y]\text{ ungerade}&[Y]\text{ gerade}\\\hline
Y'\text{ ganze gerade Zahl}&Y'\text{ ganze gerade Zahl}\\
\text{und }[X]\text{ gerade}&[X]\text{ gerade}
\end{array}
\]
\[
\begin{array}{l|l}
\frac{M}{m}\text{ im 3. Quadr.}&\frac{M}{m}\text{ im 4. Quadr.}\\
X'\text{ ganze gerade Zahl}&X'\text{ ganze gerade Zahl}\\
\text{und }[Y]\text{ gerade}&[Y]\text{ ungerade}\\\hline
Y'\text{ ganze gerade Zahl}&Y'\text{ ganze gerade Zahl}\\
\text{und }[X]\text{ ungerade}&[X]\text{ ungerade}
\end{array}
\]

\SourceLeaf{349}
So oft sich eine dieser Bedingungen in die entgegengesetzte ändert, wird $\delta=-1$; so oft sich beide ändern, bleibt $\delta=+1$.

\Article{11}
Zur bequemern Uebersicht dieser Rechnungen dienen folgende Formeln:
\[
\begin{aligned}
\text{es ist }m&=a+bi,\qquad aa+bb=d\\
M&=A+Bi,\qquad AA+BB=D\\
\frac{dM}{m}&=\alpha+\beta i,\qquad\alpha=aA+bB,\quad\beta=aB-bA\\
\frac{\xi+i\eta}{m}&=x+iy,\qquad M(x+iy)=X+iY
\end{aligned}
\]

Ist gegeben $\eta$ und $X$, so wird
\[
\begin{array}{rl}
1.&\xi=\frac{\beta\eta}{\alpha}+\frac{dX}{\alpha}\\[3pt]
2.&Y=\frac{D\eta}{\alpha}+\frac{\beta X}{\alpha}
\end{array}
\]
Ist gegeben $\eta$ und $Y$, so wird
\[
\begin{array}{rl}
3.&\xi=-\frac{\alpha\eta}{\beta}+\frac{dY}{\beta}\\[3pt]
4.&X=-\frac{D\eta}{\beta}+\frac{\alpha Y}{\beta}
\end{array}
\]
Ist gegeben $\eta$ und $x$, so wird
\[
\begin{array}{rl}
5.&\xi=-\frac{b\eta}{a}+\frac{dx}{a}\\[3pt]
6.&X=-\frac{B\eta}{a}+\frac{\alpha x}{a}\\[3pt]
7.&Y=\frac{A\eta}{a}+\frac{\beta x}{a}
\end{array}
\]
Ist gegeben $\eta$ und $y$, so wird
\[
\begin{array}{rl}
8.&\xi=\frac{a\eta}{b}-\frac{dy}{b}\\[3pt]
9.&X=\frac{A\eta}{b}-\frac{\alpha y}{b}\\[3pt]
10.&Y=\frac{B\eta}{b}-\frac{\beta y}{b}
\end{array}
\]

\Article{12}
Die Regel des $10$. Art. lässt sich nun so ausdrücken. Indem $\eta$ einen bestimmten Werth erhält, ist
\[
\Sigma\pm\Theta\frac{fM}{m}=k^0\Theta(X^0+Y^0i)-k^{00}\Theta(X^{00}+Y^{00}i)+\Sigma k
\]

\SourceLeaf{350}
Hier ist $k^0=0$, wenn $[\xi^0]$ gerade; $k^0=+1$, wenn $[\xi^0]$ ungerade und $\eta$ gerade; $k^0=-1$, wenn $[\xi^0]$ ungerade und $\eta$ ungerade ist; $k^{00}$ wird eben so durch $[\xi^{00}]$ und $\eta$ bestimmt. Endlich ist $\Sigma k$ ein Aggregat von so vielen Zahlen, als es zwischen $\xi=\xi^0$ und $\xi=\xi^{00}$ ganze Werthe von $X$ oder $Y$ gibt; jedesmal ist $k=0$, wenn das entsprechende $[\xi]$ gerade ist, hingegen $=\pm1$, wenn $[\xi]$ ungerade ist. Das Zeichen wird auf folgende Art bestimmt. Ist $X$ eine ganze Zahl, so wird $k=1$, wenn zugleich
\[
\begin{array}{ll}
\eta&\text{gerade}\\
X&\text{gerade}\\
\relax[Y]&\text{gerade}\\
\frac{M}{m}&\text{im zweiten oder dritten Quadranten d.\,i. }\alpha\text{ negativ}
\end{array}
\]

Ist eine oder drei dieser Bedingungen nicht vorhanden, so wird $k=-1$; fehlen zwei oder alle vier, so bleibt $k=1$. Ist hingegen $Y$ eine ganze Zahl, so wird $k=1$, wenn von den $4$ Bedingungen
\[
\begin{array}{ll}
\eta&\text{gerade}\\
Y&\text{gerade}\\
\relax[X]&\text{gerade}\\
\frac{M}{m}&\text{im ersten oder zweiten Quadranten d.\,i. }\beta\text{ positiv}
\end{array}
\]
alle oder zwei oder keine erfüllt ist.

\Article{13}
Jetzt haben wir noch die Fälle besonders zu betrachten, wo $\xi^0$ oder $\xi^{00}$ (oder $X^0,Y^0,X^{00},Y^{00}$) eine ganze Zahl ist. Es sind hier vier Fälle zu unterscheiden, indem wir $a$ und $A$ ungerade setzen.

I. Liegt $m$ im ersten Quadranten, so wird für $x=\frac12,y=0$; $\eta=\frac12b$ eine ganze Zahl; es ist dann $Y^{00}=\frac12B$ eine ganze Zahl und $\Theta(X^{00}+Y^{00}i)$ wird nur dann $=\Theta(\frac12A+\frac12Bi)$ sein, wenn $\beta$ negativ ist, bei einem positiven $\beta$ hingegen wird dafür $\Theta(\frac12A+(\frac12B-1)i)$ genommen werden.

II. Liegt $m$ im zweiten Quadranten, so wird für $x=0,y=0$; $\eta=0$. Hier wird für diesen Werth von $\eta$, $X^{00}=0,Y^{00}=0$. Man hat dann

\SourceLeaf{351}
\[
\begin{array}{rll}
\Theta(X^{00}+Y^{00}i)=&2,&\text{je nachdem }\dfrac{M}{m}\text{ im }1.\\
&3,&\hfill2.\\
&0&\hfill3.\\
&1&\hfill4.\text{ Quadr. liegt, und }k^{00}=1
\end{array}
\]

III. Liegt $m$ im dritten Quadranten, so wird für $x=\frac12,y=0$; $\eta=\frac12b$ eine ganze Zahl, wofür $X^0+Y^0i=\frac12A+\frac12Bi$. Man setzt dann
\[
\Theta(X^0+Y^0i)=\Theta(\tfrac12A+(\tfrac12B-1)i)
\]
so oft $\beta$ negativ ist.

IV. Liegt $m$ im vierten Quadranten, so ist für $\eta=0$,
$\Theta(X^0+Y^0i=0,1,2,3$ zu setzen, je nachdem $\frac{M}{m}$ im $1.\ 2.\ 3.\ 4.$ Quadranten liegt
$k^0=0$.

\Article{14}
Aus den vorhergehenden Untersuchungen folgt nunmehr folgende Bestimmung des Decidenten.

Man sammle alle Werthe von $x$ und $y$, die \emph{innerhalb} der Grenzen $0$ und $\frac12$ liegen und wofür entweder $\eta$ und $X$ oder $\eta$ und $Y$ eine ganze Zahl ist, und bestimme für jedes $x+iy$ nach den Regeln des $12$. Art. den Werth von $k$.

Man sammle ferner alle Werthe auf den Grenzen d.\,i. wo entweder $x=0$ oder $\frac12$, während $y$ \emph{zwischen} $0$ und $\frac12$, oder $y=0$ oder $=\frac12$, während $x$ \emph{zwischen} $0$ und $\frac12$, die so beschaffen sind, dass $\eta$ eine ganze Zahl und $[\xi]$ ungerade, und bestimme das zugehörige $l$ auf folgende Weise. Es sei $\Theta M(x+yi)=\pm\theta$, das obere Zeichen für gerade, das untere für ungerade $\eta$
\[
\begin{array}{c|c|c|c|c}
&\multicolumn{4}{c}{\text{so ist für }m\text{ im}}\\
\text{für}&1.\text{ Quadr.}&2.\text{ Quadr.}&3.\text{ Quadr.}&4.\text{ Quadr.}\\
y=0&l=-\theta&l=-\theta&l=+\theta&l=+\theta\\
x=\frac12&l=-\theta&l=+\theta&l=+\theta&l=-\theta\\
y=\frac12&l=+\theta&l=+\theta&l=-\theta&l=-\theta\\
x=0&l=+\theta&l=-\theta&l=-\theta&l=+\theta
\end{array}
\]

\SourceLeaf{352}
Kürzer so
\[
l=\pm\theta,
\]
\[
\begin{array}{ll}
\text{das Zeichen ist dasselbe wie das von }a&\text{wenn }x=0\\
\text{das entgegengesetzte}&\text{wenn }x=\frac12\\
\text{dasselbe wie das von }b&\text{wenn }y=\frac12\\
\text{das entgegengesetzte}&\text{wenn }y=0
\end{array}
\]

Zu $\Sigma k+\Sigma l$ kommt dann noch hinzu
\[
\left.\begin{array}{ll}
\text{wenn }m\text{ im zweiten Quadranten liegt:}&2,1,0,3\\
\text{wenn }m\text{ im vierten Quadranten liegt:}&0,
\end{array}\right\}
\begin{array}{l}\text{je nachdem }\frac{M}{m}\text{ im}\\1.\ 2.\ 3.\ 4.\text{ Quadr.}\end{array}
\]
wenn $m$ im ersten Quadranten liegt und $\frac12(a-1)$ ungerade ist
\[
\begin{array}{ll}
\Theta(\frac12A+(\frac12B-1)i)&\text{wenn }\beta\text{ positiv}\\
\Theta(\frac12A+\frac12Bi)&\text{wenn }\beta\text{ negativ}
\end{array}
\]
wenn $m$ im dritten Quadranten liegt und $\frac12(a-1)$ ungerade ist
\[
\begin{array}{ll}
-\Theta(\frac12A+\frac12Bi)&\text{wenn }\beta\text{ positiv}\\
-\Theta(\frac12A+(\frac12B-1)i)&\text{wenn }\beta\text{ negativ.}
\end{array}
\]
\BookRule{.3\textwidth}
